\section{Results, Deviations, and Follow-up}

\subsection{Summary of Measured Outcomes}
\begin{table}[H]
    \centering
    \small
    \begin{tabular}{p{3.2cm} p{2.1cm} p{2.1cm} p{3.8cm}}
        \toprule
        \textbf{Metric} & \textbf{Target} & \textbf{Measured} & \textbf{Interpretation} \\
        \midrule
        Startup time & $< \SI{50}{\milli\second}$ & \SI{31}{\milli\second} & Meets nominal bring-up target. \\
        Output regulation & $\pm 5\%$ & $\pm 1.8\%$ & Margin is acceptable for the current load range. \\
        Peak temperature & $< \SI{90}{\celsius}$ & \SI{74}{\celsius} & Thermal headroom remains available. \\
        Fault response latency & $< \SI{5}{\milli\second}$ & Pending & Requires high-speed trigger validation. \\
        \bottomrule
    \end{tabular}
    \caption{Representative measured outcomes}
    \label{tab:measured-outcomes}
\end{table}

\subsection{Observed Deviations}
\riskbox{
List every blocked or partially executed test here. A deferred case is not a
minor note; it changes the decision quality of the report.
}

\makeactiontable{Open Actions After Test Campaign}{
    Repeat abnormal-load protection capture with approved current-limit settings. & Safety Lead & 2026-04-15 & Open \\
    Replace placeholder plots with signed scope exports from the final bench run. & Test Eng. & 2026-04-12 & In Progress \\
    Review trigger configuration for the fault-latency measurement. & Systems Eng. & 2026-04-11 & Open \\
}

\subsection{Conclusion}
\decisionbox{
The document should end with a direct statement: either the test objective was
met, met with restrictions, or remains open. Avoid vague language when the
report is part of a release or approval gate.
}
